\section{Common Mode and Centered Relational Contrast}

\subsection{Orthogonal projections}

Let
\[
J=\one_{15}\one_{15}^{\mathsf T}.
\]
Define
\[
\Pzero
=
\frac{1}{15}J
\]
and
\[
\Pperp
=
I-\Pzero.
\]

Then \(\Pzero\) is the orthogonal projection onto
\[
\operatorname{span}\{\one_{15}\},
\]
and \(\Pperp\) is the orthogonal projection onto
\[
\one_{15}^{\perp}
=
\left\{
x\in\mathbb R^{15}:
\one_{15}^{\mathsf T}x=0
\right\}.
\]

Every \(x\in\mathbb R^{15}\) has the unique decomposition
\[
x
=
\Pzero x+\Pperp x.
\]

Writing
\[
\overline{x}
=
\frac{1}{15}\one_{15}^{\mathsf T}x,
\]
we have
\[
\Pzero x
=
\overline{x}\,\one_{15}.
\]

\subsection{Common and relational magnitudes}

Define
\[
\Ccommon(x)
:=
\|\Pzero x\|_2
=
\sqrt{15}\,|\overline{x}|
\]
and
\[
\Rrel(x)
:=
\|\Pperp x\|_2
=
\left\|
x-\overline{x}\,\one_{15}
\right\|_2.
\]

The second quantity is exactly the centered relational magnitude
specialized to the fifteen-sector register.

Orthogonality gives
\[
\|x\|_2^2
=
\Ccommon(x)^2+\Rrel(x)^2.
\]

\subsection{Dimensionless contrast fraction}

For \(x\neq0\), define
\[
\zetaRel(x)
:=
\frac{\Rrel(x)^2}{\|x\|_2^2}.
\]
Then
\[
0\le\zetaRel(x)\le1.
\]

Moreover,
\[
\zetaRel(x)=0
\quad\Longleftrightarrow\quad
x\in\operatorname{span}\{\one_{15}\},
\]
while
\[
\zetaRel(x)=1
\quad\Longleftrightarrow\quad
x\in\one_{15}^{\perp}.
\]

Thus \(\zetaRel\) measures the fraction of Euclidean register
magnitude carried by centered relational contrast rather than by the
common mode.

When \(\Ccommon(x)\neq0\), we also define the relative contrast ratio
\[
\rho(x)
:=
\frac{\Rrel(x)}{\Ccommon(x)}.
\]

\subsection{Invariant decomposition under the quadratic}

Because \(Q\) is a polynomial in the regular graph adjacency matrix
\(A\), both \(\Pzero\) and \(\Pperp\) commute with \(Q\), and hence
with \(\widehat{Q}\):
\[
\Pzero\widehat{Q}
=
\widehat{Q}\Pzero,
\qquad
\Pperp\widehat{Q}
=
\widehat{Q}\Pperp.
\]

Therefore the common and centered subspaces are invariant.

This point is essential.  Contrast concentration under
\(\widehat{Q}\) does not arise because centered amplitude is
transferred into the common mode.  The two subspaces do not mix.
Instead, the common mode has eigenvalue \(1\), while every centered
mode has eigenvalue strictly smaller than \(1\).

\subsection{Complete spectral decomposition}

Let \(E_\lambda\) denote the orthogonal eigenspace of \(A\) with
eigenvalue \(\lambda\).  Then
\[
\mathbb R^{15}
=
E_4
\oplus
E_2
\oplus
E_{-1}
\oplus
E_{-2},
\]
with
\[
\dim E_4=1,\quad
\dim E_2=5,\quad
\dim E_{-1}=4,\quad
\dim E_{-2}=5.
\]

The common-mode space is
\[
E_4=\operatorname{span}\{\one_{15}\}.
\]

On these four spaces, \(\widehat{Q}\) acts respectively by
\[
1,\qquad
\frac{9}{49},\qquad
\frac{3}{98},\qquad
\frac{1}{49}.
\]

The largest centered eigenvalue is therefore
\[
\lambda_{\mathrm{cent}}
=
\frac{9}{49}.
\]
